%-------------------------------------
% 中文简历 — 一页版 (XeLaTeX)
% 编译: xelatex -interaction=nonstopmode -halt-on-error cv_zh_1page.tex
% 字体: 默认使用 fandol 字体 (随 TeX 发行版自带, 跨平台均可编译)。
%       若在 Windows 上偏好系统字体 (宋体 SimSun / 黑体 SimHei), 将 fontset=fandol 改为 fontset=windows。
%-------------------------------------
\documentclass[a4paper,9pt]{article}

\usepackage[UTF8, fontset=fandol, scheme=plain]{ctex}
\usepackage[a4paper,left=1.4cm,right=1.4cm,top=1.0cm,bottom=1.0cm]{geometry}
\usepackage{myresume_cjk}
\usepackage[unicode]{hyperref}
\hypersetup{colorlinks=true, urlcolor=Accent, linkcolor=Accent}

\linespread{0.98}
% 一页版: 适当压缩段落间距
\renewcommand{\cvsection}[1]{%
  \par\vspace{5pt}%
  {\heiti\fontsize{12pt}{13pt}\selectfont\color{Accent}#1}\par
  \vspace{1.5pt}{\color{RuleColor}\hrule height 0.6pt}\vspace{3.5pt}%
}

\begin{document}

% -------------------- 抬头 --------------------
\begin{minipage}[t]{0.58\textwidth}
  \ResumeName{王瑞文}\;\ResumeNameEn{Ruiwen Wang}\\[3pt]
  \ResumeRole{AI 基础设施 \textperiodcentered{} 大模型预训练 \textperiodcentered{} 混合并行}
\end{minipage}
\hfill
\begin{minipage}[t]{0.40\textwidth}
  \raggedleft\footnotesize\color{TextMuted}
  法国，巴黎 \;|\; (+33) 663113466\\
  \href{mailto:wangrw0124@gmail.com}{wangrw0124@gmail.com} \;|\; 国籍：法国\\
  \href{http://ruiwen.wang/}{ruiwen.wang} \textperiodcentered\ \href{https://scholar.google.com/citations?user=AR1nHEUAAAAJ}{Scholar} \textperiodcentered\ \href{https://github.com/wang-ruiwen}{GitHub} \textperiodcentered\ \href{https://www.linkedin.com/in/ruiwen-wang-ai/}{LinkedIn}
\end{minipage}

\vspace{6pt}
\AccentRule

% -------------------- 个人简介 --------------------
\cvsection{个人简介}
\footnotesize
计算机科学博士研究生（CIFRE 产学联合培养），就读于索邦大学与 EURECOM，并与华为巴黎研究所联合开展研究，预计 2026 年 8 月毕业。研究位于 AI、高性能计算（HPC）与机器学习系统（MLSys）交叉领域，聚焦超大规模 LLM/DNN 训练的优化式策略规划：以符号化性能 / 显存模型结合 ILP 与约束求解，自动生成显存可行、通信感知的混合并行策略与流水线调度。成果发表于 CLUSTER、Euro-Par（最佳海报奖）、NPC、HPC Asia，并在最多 16,384 张 Ascend NPU 的生产集群上验证落地。

% -------------------- 工作经历 --------------------
\cvsection{工作经历}
\cventry
  {华为巴黎研究所 \textemdash{} 分布式与并行计算实验室}{法国，巴黎}
  {博士研究生（CIFRE）兼研究工程师 \textemdash{} 大模型训练系统}{2021 \textemdash{} 2026（预计）}
\itemListStart
  \myItem{构建基于代价模型的混合 / 流水线并行 LLM 训练规划器：显存与性能约束下的阶段划分、流水线调度与重计算选择。}
  \myItem{设计符号化峰值显存估计与集合通信代价模型，结合通信-计算重叠策略，进行大规模吞吐与 MFU 调优。}
  \myItem{将研究原型产品化为可部署规划流程，与工程及客户团队协作落地；推动研究向生产级训练环境的技术转化。}
\itemListEnd
{\footnotesize\color{TextMuted}角色演进：研究学徒（2021\textendash2022）$\rightarrow$ 研究工程师（2022\textendash2023）$\rightarrow$ CIFRE 博士研究生（2023\textendash2026）。}

% -------------------- 核心成果 --------------------
\cvsection{核心成果}
\itemListStart
  \myItem{\textbf{BMPipe}（CLUSTER 2025）：符号代价模型 + ILP 联合优化阶段边界 / 分块 / 重计算；最高 1.36$\times$ 加速，HBM 利用率 83\textendash92\%，验证至 16{,}384 张 Ascend NPU。}
  \myItem{\textbf{ManuMatic}（NPC 2025）：算子切分作为一等规划约束；相较 D-Rec 在 Mixtral-8$\times$7B 上 2.24$\times$、Llama3-8B 上 2.04$\times$，Qwen2.5-72B 上较专家方案 +30.1\%。}
  \myItem{\textbf{H2O}（Euro-Par 2025 最佳海报奖）：两级规划工作流；MFU 26.13\%$\rightarrow$35.73\%，DeepSeek-141B 加速 36.72\%，并转生产（4{,}096 卡约 1.1$\times$，512 卡长序列最高 2.37$\times$）。}
  \myItem{\textbf{PRISM}（HPC Asia 2026）：免性能剖析的符号化规划器；1.23$\times$\textendash1.43$\times$ 加速，MFU 29.40\%$\rightarrow$36.20\%，验证至 1{,}024 卡。}
\itemListEnd

% -------------------- 论文与专利 --------------------
\cvsection{论文与专利}
\footnotesize
\textbf{论文（第一/共同作者，节选）}：BMPipe（IEEE CLUSTER 2025）；H2O（Euro-Par 2025，最佳海报奖）；ManuMatic（NPC 2025）；Training Report of TeleChat3-MoE（arXiv 2025）；PRISM 与 MLLM Pipeline Bubble Modeling（HPC Asia 2026）。\par
\vspace{2pt}
\textbf{专利}：3 项 —— 机器学习模型执行计划生成（WO 2025/020165 A1）、大模型负载均衡方法、大模型训练多维并行配置方法。

% -------------------- 教育背景 --------------------
\cvsection{教育背景}
\footnotesize
\textbf{索邦大学 \& EURECOM} —— 计算机科学博士（导师：Prof. Raja Appuswamy、Prof. Chong Li）\hfill {\color{TextMuted}2023 \textemdash{} 2026（预计）}\\
\textbf{索邦大学} —— 计算机科学硕士（算法、软件科学与技术）\hfill {\color{TextMuted}2019 \textemdash{} 2022}\\
\textbf{巴黎-萨克雷大学} —— 计算机科学学士 \hfill {\color{TextMuted}2016 \textemdash{} 2019}

% -------------------- 荣誉 与 技能 --------------------
\begin{minipage}[t]{0.48\textwidth}
\cvsection{荣誉奖项}
\awardListStart
  \item \footnotesize 金奖，华为自动负载均衡挑战赛 \hfill 2024.08
  \item \footnotesize 最佳海报奖，Euro-Par 国际会议 \hfill 2025.08
  \item \footnotesize 感谢信，中国电信（DeepSeek-MoE 部署）\hfill 2025.12
\awardListEnd
\end{minipage}
\hfill
\begin{minipage}[t]{0.48\textwidth}
\cvsection{技能}
\footnotesize
\textbf{编程}：Python、C++、Java\\
\textbf{框架/平台}：MindSpore、PyTorch；大规模 GPU/NPU 集群与云\\
\textbf{工具/开发}：Ascend 性能剖析、tracing、性能基准；Linux、Git\\
\textbf{语言}：中文（母语）、英语、法语、粤语
\end{minipage}

\end{document}
